\Needspace{18\baselineskip}
\subsection{Arbitrary curve backgrounds}
\label{subsec:arbitrary-curve-product-metrics}

Scaling by coefficients is a special case of changing the curve backgrounds. By \cref{prop:separated-curve-obstacle-projections}, the coordinatewise computation of the envelopes holds for arbitrary smooth curve backgrounds of nonnegative area, so the product metric depends on the areas of the curve backgrounds rather than on their positivity. A curve background of area zero has a unique normalized potential, which is smooth. Its coordinate no longer contributes to the metric, but its smooth potential remains a summand of the product potentials.

Keep $M=Y\times C_0\times Z$, the projections $p,s,r$ and $\Psi$ from \cref{cor:nonclosed-separated-product-metric}. On each $C_i$, choose a smooth real $(1,1)$-form $\sigma_i$ with $\PSH(C_i,\sigma_i)\ne\varnothing$, and put $B_i=\int_{C_i}\sigma_i\geq0$. For a smooth function $f\geq0$ on $Y$, set
\[
 \gamma_\sigma=\sum_i\pi_i^*\sigma_i,\qquad
 \Theta_{\sigma,f}=p^*\gamma_\sigma+(f\circ p)s^*\tau+r^*\kappa.
\]
For normalized ordinary models $q_i$ in the curve backgrounds $\sigma_i$, write
\[
 \Phi_\sigma(q)=\Psi+\sum_i q_i\circ\pi_i\circ p,
 \qquad q_{i,\mathcal I}=P_{\sigma_i}[q_i]_{\mathcal I},
\]
and let $\mathcal S_\sigma$ be the locus of these potentials. All coordinates are included: when $B_i=0$, $q_i$ is the unique normalized potential $u_{\sigma_i}$ of \cref{lem:curve-background-area}, which is smooth.

\Needspace{14\baselineskip}
\begin{lemma}\label{lem:arbitrary-product-model-projections}\index{Envelope!arbitrary curve product backgrounds}
Every $\Phi_\sigma(q)$ is a normalized ordinary model, and
\[
 P_{\Theta_{\sigma,f}}[\Phi_\sigma(q)]_{\mathcal I}
       =\Phi_\sigma(q_{\mathcal I}).
\]
The background is not big if $f=0$ or some $B_i=0$. If $f>0$ everywhere and all $B_i>0$, it is big. It is not closed whenever $f$ is nonconstant.
\end{lemma}
\begin{proof}
On $B=Y\times C_0$, let $\pi,s_0$ be the projections and put
\[
 \gamma_{\sigma,f}=\pi^*\gamma_\sigma+(f\circ\pi)s_0^*\tau,
 \qquad w=\sum_iq_i\circ\pi_i.
\]
The case of the zero obstacle in \cref{prop:separated-curve-obstacle-projections} gives $P_{\gamma_\sigma}[w]=w$. Replace an ordinary competitor for $w\circ\pi$ by its maximum with $w\circ\pi$. Its restriction to every section $Y\times\{z\}$ is then a competitor for $P_{\gamma_\sigma}[w]$ that is not identically $-\infty$, so
\[
 P_{\gamma_{\sigma,f}}[w\circ\pi]=w\circ\pi.
\]

Use the atomic Green representatives $g_i$ constructed in the proof of \cref{prop:separated-curve-obstacle-projections}, which are smooth in the coordinates of zero area. Their sum $g=\sum_i g_i\circ\pi_i\circ\pi$ is $\gamma_{\sigma,f}$-plurisubharmonic and multiplier-ideal equivalent to $w\circ\pi$, and its complex Hessian is a closed countable divisor current of finite mass minus a smooth form. If $w_{\mathcal I}=\sum_iq_{i,\mathcal I}\circ\pi_i$, then $g\sim w_{\mathcal I}\circ\pi$. By \cref{lem:divisorial-potential-domination}, every ideal competitor is ordinarily dominated by $g$, hence bounded by
\[
 P_{\gamma_{\sigma,f}}[g]
    =P_{\gamma_{\sigma,f}}[w_{\mathcal I}\circ\pi]
    =w_{\mathcal I}\circ\pi.
\]
Here the last equality holds because each ideal projection $q_{i,\mathcal I}$ on a curve is an ordinary model, so the argument with sections just given applies to $w_{\mathcal I}$. By pullback and addition, $w_{\mathcal I}\circ\pi$ is itself an ideal competitor, which proves the ideal identity on $B$. Now apply \cref{prop:fibration-obstacle-reduction} to $M\to B$ and the divisor $B\times D$ with weight one, whose induced zero obstacle is zero. The two identities upstairs and the independence of the factors give the assertion and normalization.

For the statements on bigness, put $n=m+1+k$ and $\Omega=p^*\beta+s^*\tau+r^*\kappa>0$, with the positive product form $\beta$ fixed earlier. For $f=0$, the closed positive test form
\[
 T_0=p^*\left(\bigwedge_j\pi_j^*\eta_j\right)\wedge r^*(\kappa^k)
\]
satisfies $\int_M\Theta_{\sigma,0}\wedge T_0=0$ and $\int_M\Omega\wedge T_0>0$. If $B_i=0$, use instead
\[
 T_i=p^*\left(\bigwedge_{j\ne i}\pi_j^*\eta_j\right)
             \wedge s^*\tau\wedge r^*(\kappa^k).
\]
Then $\int_M\Theta_{\sigma,f}\wedge T_i=0$ for every $f$, since the only nonzero term contains the factor $\int_{C_i}\sigma_i=B_i$, while $\int_M\Omega\wedge T_i>0$. In either case, Stokes' theorem against the closed test form shows that no current $\Theta_{\sigma,f}+\ddc u$ with $u$ quasi-plurisubharmonic is bounded below by a positive multiple of $\Omega$, so the background is not big.

Finally, write $\sigma_i=B_i\eta_i/A_i+\ddc\chi_i$ by \cref{lem:curve-background-area}. The smooth potential $-\sum_i\chi_i\circ\pi_i\circ p$ has strictly positive current when every $B_i>0$ and $f>0$, which proves bigness in that case. The identity $\mathrm d\Theta_{\sigma,f}=p^*(\mathrm df)\wedge s^*\tau$ shows that the background is not closed for nonconstant $f$.
\end{proof}

For the polar curvatures $\mu_i$ of $q_i$ and $\nu_i$ of another tuple $\widetilde q$, taken in the backgrounds $\sigma_i$, set
\begin{equation}\label{eq:arbitrary-product-polar-metric}
 \rho_\sigma(\Phi_\sigma(q),\Phi_\sigma(\widetilde q))
       =\sum_i\|\mu_i-\nu_i\|_{\mathrm{TV}}.
\end{equation}
Write $\mathcal I_\sigma\subset\mathcal S_\sigma$ for the ideal subspace and $\mu_i^{\mathrm{na}}$ for the nonatomic part of $\mu_i$.

\Needspace{20\baselineskip}
\begin{proposition}\label{prop:arbitrary-curve-product-metric}\index{Metric completeness!arbitrary curve product backgrounds}\index{Model potential!area comparison}
The metric space $(\mathcal S_\sigma,\rho_\sigma)$ is complete and geodesic, with diameter $2\sum_i B_i$. The ideal subspace is closed, geodesically convex and equal to the metric closure of the models with analytic singularities. Ideal projection is the unique nearest-point retraction onto it and is $1$-Lipschitz, and the distance to the subspace is $\sum_i\mu_i^{\mathrm{na}}(C_i)$. Metric convergence implies $L^1(M)$ convergence for every smooth positive volume form. These assertions are independent of $f$, and the upper-volume comparison vanishes on this locus.

Let $\sigma'_i$ be another tuple of smooth real forms with $\PSH(C_i,\sigma'_i)\ne\varnothing$, with areas $B'_i$. Then
\begin{equation}\label{eq:arbitrary-product-area-comparison}
 \left|\sum_i(B_i-B'_i)\right|
 \leq d_{\mathrm{GH}}(\mathcal S_\sigma,\mathcal S_{\sigma'})
 \leq\sum_i|B_i-B'_i|.
\end{equation}
The same bounds hold for the ideal subspaces. If $B'_i\leq B_i$ for every $i$, both distances equal $\sum_i(B_i-B'_i)$. In this case, rescaling the polar measures by the area ratios, with the unique normalized potential in each target coordinate of zero area, gives a surjective $1$-Lipschitz map that commutes with ideal projection. Its fibers have diameter $2\sum_{\{i:B_i>0,\ B'_i=0\}}B_i$.

Both $\mathcal S_\sigma$ and $\mathcal I_\sigma$ are noncompact if some $B_i>0$. If all $B_i=0$, they coincide with the singleton $\{\Psi+\sum_i u_{\sigma_i}\circ\pi_i\circ p\}$.
\end{proposition}
\begin{proof}
Restriction to coordinate sections on which the other potentials are finite determines each $q_i$ up to a constant, and normalization fixes that constant. The preceding lemma and \cref{prop:curve-arbitrary-background-models} therefore identify the locus with the finite product of the spaces of curve models, with the sum of their polar metrics; factors of zero area are single points. Completeness, geodesics and the diameter formula follow coordinatewise. Under this identification, \cref{cor:curve-model-completeness,prop:curve-analytic-distance-defect} give closedness, geodesic convexity and the unique nearest-point retraction onto the ideal subspace, and the distances add. Fiber integration transfers $L^1$ convergence on the curves to $M$.

For analytic density, approximate each atomic polar measure by finite atomic measures with rational coefficients and mass at most $B_i$. By the atomic Green construction and transport under smooth gauge changes, their models in background $\sigma_i$ have the corresponding finitely many poles with rational coefficients, up to bounded remainders. Coordinates of zero area are smooth. Together with the unit coefficient of $D$, the poles in the coordinates of positive area allow the denominators to be cleared in local product monomials, which gives analytic models in the locus. They converge for the metric \eqref{eq:arbitrary-product-polar-metric}. Conversely, by \cref{prop:analytic-type-calculus}, a product model with analytic singularities restricts, on every coordinate curve section on which the other potentials are finite and whose $Z$-coordinate lies outside $D$, to a function with analytic singularities that is not identically $-\infty$. Each $q_i$ therefore has polar curvature with finitely many atoms, and it is an ideal model by the classification of curve models. Closedness proves the analytic-closure assertion.

By \cref{prop:fibration-obstacle-reduction}, applied to the divisor $B\times D$ with weight one, every model and every pairwise maximum in this locus has zero upper volume; for a maximum, use $\Psi$ plus the pullback of the maximum of the two potentials on $B$. Thus all three terms of the upper-volume comparison vanish, even in the backgrounds that are not big by the lemma.

To compare different backgrounds, put $a_i=B_i/A_i$, and for a positive polar measure $\mu$ on $C_i$ of mass at most $B_i$ write $p_{\sigma_i}(\mu)$ for the ordinary model in background $\sigma_i$ with polar curvature $\mu$, as in \cref{prop:curve-arbitrary-background-models}. Sending a curve model for $\eta_i$ with polar measure $\mu_i$, of mass at most $A_i$, to $p_{\sigma_i}(a_i\mu_i)$, which for $a_i=0$ is the unique normalized potential, defines an isometry from the weighted locus $\mathcal S_a$ of \cref{prop:weighted-product-polar-metric} onto $\mathcal S_\sigma$ that maps the ideal subspaces onto each other. Applying \cref{prop:weighted-product-metric-comparison} with $a_i=B_i/A_i$ and $b_i=B'_i/A_i$ gives \eqref{eq:arbitrary-product-area-comparison}, the exact value for ordered areas, the contractions and the diameters of their fibers. In curve coordinates the contraction is
\[
 p_{\sigma_i}(\mu)\longmapsto
 \begin{cases}
 p_{\sigma'_i}((B'_i/B_i)\mu),&B_i>0,\\
 p_{\sigma'_i}(0),&B_i=0.
 \end{cases}
\]
In the second case $B'_i=0$, because $B'_i\leq B_i$. Atomic projection commutes with both operations. As in \cref{thm:curve-background-homothety}, irrational area ratios need not preserve individual analytic models, but the ideal subspaces are preserved for every ratio.

If $B_i>0$, the models with polar measure $B_i\delta_x$ for infinitely many distinct points $x\in C_i$, with all other coordinates fixed at their models with zero polar curvature, form an infinite set of ideal product models with mutual distances $2B_i$. This excludes total boundedness of both loci. If all areas are zero, every curve coordinate is fixed at its unique normalized smooth potential, which proves the last assertion.
\end{proof}

\Needspace{16\baselineskip}
\begin{corollary}\label{cor:arbitrary-product-zero-area-collapse}\index{Model potential!smooth zero-area product limit}
For each $i$, let $\sigma_{i,j}$ be smooth real forms on $C_i$ of nonnegative area that converge uniformly, as $j\to\infty$, to a smooth form $\sigma_{i,0}$ of area zero, and write $\sigma_j=(\sigma_{1,j},\ldots,\sigma_{m,j})$. For any nonnegative smooth auxiliary coefficients $f_j$, the diameters of $\mathcal S_{\sigma_j}$ tend to zero. For every fixed smooth positive volume form on $M$, the potentials in $\mathcal S_{\sigma_j}$ converge in $L^1(M)$ to
\[
 \Psi+\sum_i u_{\sigma_{i,0}}\circ\pi_i\circ p,
\]
uniformly over all normalized model tuples.
\end{corollary}
\begin{proof}
The areas $B_{i,j}$ tend to zero, so the diameter formula applies. By \cref{thm:curve-zero-area-collapse}, all normalized $\sigma_{i,j}$-plurisubharmonic potentials converge in $L^1(C_i)$ to $u_{\sigma_{i,0}}$, uniformly over this family. Integration along the fixed compact fibers bounds each pulled-back norm by a constant times the norm on the curve, and the finite sum gives the asserted uniform limit on $M$. The auxiliary coefficients do not enter these potentials or their metrics.
\end{proof}

Thus the metric comparison sees only the areas of the curve backgrounds, while the limiting potential still records their smooth gauges. These statements concern the separated product locus; extending the metric to general ordinary models is part of \cref{qu:i-good-no-bmp}.
